\begingroup
\renewcommand{\arraystretch}{1.45}
\fancyhead[L]{\small BESTIMMUNG DER ANZIEHUNG}
\Needspace{37\baselineskip}
\AttractionArticle{4}

Setzt man jetzt im Ausdruck für den Abstand $\rho$ für $\cos E$ und $\sin E$ die oben angenommenen Werte ein, so wird er in die folgende Form übergehen:
\[
\rho=\frac{\surd\left(\begin{gathered}G+G'\cos T^2+G''\sin T^2\\[5pt]
+2H\cos T\sin T+2H'\sin T+2H''\cos T\end{gathered}\right)}{\gamma+\gamma'\cos T+\gamma''\sin T}
\]
Dabei werden wir die Koeffizienten $\alpha$, $\alpha'$, $\alpha''$ usw. so bestimmen, dass unter Wahrung der sechs Bedingungsgleichungen
\[
\left.\begin{aligned}
-\alpha\alpha-\beta\beta+\gamma\gamma&=1\\[5pt]
-\alpha'\alpha'-\beta'\beta'+\gamma'\gamma'&=-1\\[5pt]
-\alpha''\alpha''-\beta''\beta''+\gamma''\gamma''&=-1\\[5pt]
-\alpha'\alpha''-\beta'\beta''+\gamma'\gamma''&=0\\[5pt]
-\alpha''\alpha-\beta''\beta+\gamma''\gamma&=0\\[5pt]
-\alpha\alpha'-\beta\beta'+\gamma\gamma'&=0
\end{aligned}\right\}\quad\text{[1]}
\]
und damit auch der übrigen aus ihnen hervorgehenden gilt
\[
\begin{aligned}
&H=0,\quad H'=0,\quad H''=0
\end{aligned}
\]
\Needspace{24\baselineskip}
Auf diese Weise wird das Problem, allgemein gesprochen, bestimmt sein. Bezeichnen wir also den Nenner von $\rho$ mit $t$, so muss die folgende Funktion der drei Größen $t$, $t\cos E$, $t\sin E$
\[
\begin{aligned}
&(AA+BB+CC)tt+aa(t\cos E)^2+bb(t\sin E)^2-2aAt\, .\, t\cos E\\[6pt]
&\qquad{}-2bBt\, .\, t\sin E
\end{aligned}
\]
durch die Einsetzung
\[
\begin{aligned}
&t\cos E=\alpha+\alpha'\cos T+\alpha''\sin T\\[6pt]
&t\sin E=\beta+\beta'\cos T+\beta''\sin T\\[6pt]
&t=\gamma+\gamma'\cos T+\gamma''\sin T
\end{aligned}
\]
\SourceLeaf{348}
in
\[
\begin{aligned}
&G+G'\cos T^2+G''\sin T^2
\end{aligned}
\]
übergehen. \Needspace{28\baselineskip}
Offenbar ist dies dasselbe, als wenn du sagtest, die folgende Funktion der drei Unbestimmten $x,y,z$, nämlich $(W)$,
\[
\begin{aligned}
&aaxx+bbyy+(AA+BB+CC)zz-2aAxz-2bByz
\end{aligned}
\]
müsse durch die Einsetzung
\[
\begin{aligned}
&x=\alpha u+\alpha'u'+\alpha''u''\\[6pt]
&y=\beta u+\beta'u'+\beta''u''\\[6pt]
&z=\gamma u+\gamma'u'+\gamma''u''
\end{aligned}
\]
in die folgende Funktion der Unbestimmten $u,u',u''$
\[
\begin{aligned}
&Guu+G'u'u'+G''u''u''
\end{aligned}
\]
übergehen. \Needspace{36\baselineskip}
Da aber aus diesen Formeln mithilfe der Gleichungen [1] leicht folgt
\[
\begin{aligned}
&u=-\alpha x-\beta y+\gamma z\\[6pt]
&u'=\alpha'x+\beta'y-\gamma'z\\[6pt]
&u''=\alpha''x+\beta''y-\gamma''z
\end{aligned}
\]
wird die Funktion $W$ offenbar mit der folgenden identisch sein müssen
\[
\begin{aligned}
&G(-\alpha x-\beta y+\gamma z)^2
+G'(\alpha'x+\beta'y-\gamma'z)^2
\\[6pt]
&\qquad{}+G''(\alpha''x+\beta''y-\gamma''z)^2
\end{aligned}
\]
woraus wir sechs Gleichungen erhalten
\[
\left.
\begin{aligned}
aa&=G\alpha\alpha+G'\alpha'\alpha'+G''\alpha''\alpha''\\[5pt]
bb&=G\beta\beta+G'\beta'\beta'+G''\beta''\beta''\\[5pt]
AA+BB+CC&=G\gamma\gamma+G'\gamma'\gamma'+G''\gamma''\gamma''\\[5pt]
bB&=G\beta\gamma+G'\beta'\gamma'+G''\beta''\gamma''\\[5pt]
aA&=G\gamma\alpha+G'\gamma'\alpha'+G''\gamma''\alpha''\\[5pt]
0&=G\alpha\beta+G'\alpha'\beta'+G''\alpha''\beta''
\end{aligned}
\right\}\quad [2]
\]

Aus diesen zwölf Gleichungen [1] und [2] werden unsere Unbekannten $G,\allowbreak G',\allowbreak G'',\allowbreak \alpha,\allowbreak \alpha',\allowbreak \alpha''$ usw. zu bestimmen sein.
\SourceLeaf{349}
\Needspace{36\baselineskip}
\AttractionArticle{5}

Durch Verbindung der Gleichungen [1] und [2] werden die folgenden leicht hergeleitet:
\[
\begin{aligned}
&-\alpha aa+\gamma aA=\alpha G\\[6pt]
&-\beta bb+\gamma bB=\beta G\\[6pt]
&\gamma(AA+BB+CC)-\alpha aA-\beta bB=\gamma G
\end{aligned}
\]
woraus weiter folgt
\[
\begin{aligned}
&\alpha=\frac{\gamma aA}{aa+G}\quad\mathinner{.\,.\,.\,.\,.\,.\,.\,.\,.\,.\,.\,.\,.\,.\,.\,.\,.}\quad\text{[3]}\\[6pt]
&\beta=\frac{\gamma bB}{bb+G}\quad\mathinner{.\,.\,.\,.\,.\,.\,.\,.\,.\,.\quad.\,.\,.\,.\,.\,.}\quad\text{[4]}\\[6pt]
&AA+BB+CC-\frac{aaAA}{aa+G}-\frac{bbBB}{bb+G}=G
\end{aligned}
\]
Die letzte können wir auch so ausdrücken
\[
\begin{aligned}
&\frac{AA}{aa+G}+\frac{BB}{bb+G}+\frac{CC}{G}=1\quad\mathinner{.\,.\,.\,.\,.\,.}\quad\text{[5]}
\end{aligned}
\]

\Needspace{37\baselineskip}
Ebenso leiten wir durch Verbindung der Gleichungen [1] und [2] her
\[
\begin{aligned}
&\alpha'aa-\gamma'aA=\alpha'G'\\[6pt]
&\beta'bb-\gamma'bB=\beta'G'\\[6pt]
&-\gamma'(AA+BB+CC)+\alpha'aA+\beta'bB=\gamma'G'
\end{aligned}
\]
und hieraus
\[
\begin{aligned}
&\alpha'{}=\frac{\gamma'aA}{aa-G'}\quad\mathinner{.\,.\,.\,.\,.\,.\,.\,.\,.\,.\,.\,.\,.\,.\,.\,,\,.\,.}\quad\text{[6]}\\[6pt]
&\beta'{}=\frac{\gamma'bB}{bb-G'}\quad\mathinner{.\,.\,.\,.\,.\,.\,.\,.\,.\,.\,.\,.\,.\,.\,.\,.\,.\,.}\quad\text{[7]}\\[6pt]
&\frac{AA}{aa-G'}+\frac{BB}{bb-G'}-\frac{CC}{G'}=1\quad\mathinner{.\,.\,.\,.\quad.\,.}\quad\text{[8]}
\end{aligned}
\]
\Needspace{34\baselineskip}
und auf ganz ähnliche Weise
\[
\begin{aligned}
&\alpha''{}=\frac{\gamma''aA}{aa-G''}\quad\mathinner{.\,.\,.\,.\,.\,.\,.\,.\,.\,.\,.\,.\,.\,.\,.\,.\,.\,.}\quad\text{[9]}\\[6pt]
&\beta''{}=\frac{\gamma''bB}{bb-G''}\quad\mathinner{.\,.\,.\,.\,.\,.\,.\,.\,.\,.\,.\,.\,.\,.\,.\,.\,.}\quad\text{[10]}\\[6pt]
&\frac{AA}{aa-G''}+\frac{BB}{bb-G''}-\frac{CC}{G''}=1\quad\mathinner{.\,.\,.\,.\,.\,.}\quad\text{[11]}
\end{aligned}
\]
Somit ist offenbar, dass $G,-G',-G''$ die Wurzeln der Gleichung sind
\[
\begin{aligned}
&\frac{AA}{aa+x}+\frac{BB}{bb+x}+\frac{CC}{x}=1\quad\mathinner{.\,.\,.\,.\,.\,.\,.\,.}\quad\text{[12]}
\end{aligned}
\]
\SourceLeaf{350}
die, gehörig entwickelt, folgendermaßen lautet
\[
\begin{aligned}
&x^3-(AA+BB+CC-aa-bb)xx\\[6pt]
&\qquad{}+(aabb-aaBB-aaCC-bbAA-bbCC)x\\[6pt]
&-aabbCC=0\quad\mathinner{.\,.\,.\,.\,.\,.\,.\,.\,.\,.\,.\,.\,.\,.\,.\,.}\quad\text{[13]}
\end{aligned}
\]

\clearpage
\endgroup
